% !TeX root = ../main.tex
\begin{abstract}
\label{abstract}
Sequential Recommendation (SR) learns sequence-level Collaborative Filtering (CF) from interactions, but ID-based models struggle with sparse data and cold-start items.
In contrast, item titles and images can reveal preference continuity, such as a recurring color or style.
Yet existing approaches miss visual continuity: small encoders take a single item as input, and Large Language Models (LLMs) read only titles.
We therefore shift our focus to Large Vision-Language Models (VLMs), which can read a user's entire image sequence, not just its titles.
However, VLMs inherently suffer from \textbf{modality collapse}: they over-rely on text and, as a result, overlook visual continuity for recommendation.
We find that standard contrastive Supervised Fine-Tuning (SFT), essential for adapting VLMs to SR, further amplifies this collapse: updates follow the text gradient, visual discrimination weakens, and adding image embeddings to the text embeddings consistently lowers accuracy.
We propose \textbf{VLM2Rec} with two objectives: \textit{Weak-modality Penalized Contrastive Learning} strengthens each user's weaker modality, and \textit{Cross-modal Relational Topology Regularization} aligns each user's relative similarity structure over the same items across the two modality spaces.
Its variant, $\text{VLM2Rec}_{\text{plug}}$, reuses the learned item embeddings in SR backbones and thus serves at the latency of conventional SR models.
Experiments show that VLM2Rec mitigates the collapse and consistently improves accuracy and robustness.
% Our code is available at {\color{blue}\url{https://anonymous.4open.science/r/VLM2Rec/}}.
\end{abstract}
